\chapter{Methodology}
\label{cha:methodology}

The observed-data likelihood under the model of Chapter~\ref{cha:model-specification} is

\begin{equation}
    L(\bm{\theta})=\prod_{j=1}^{J}\prod_{i\in \mathcal{I}_{j}} p(\bm{y}_{ij}\mid \bm{\theta})=\prod_{j=1}^{J}\prod_{i\in \mathcal{I}_{j}} \int p(\bm{y}_{ij}\mid \bm{\lambda}_{j})p(\bm{\lambda}_{j}\mid \bm{\Sigma})\,d\bm{\lambda}_{j},
    \label{eq:log-lik-chap2}
\end{equation}
where the parameter space is $\bm{\theta}:=\{\bm{\mu}, \bm{\Phi}, \mathbf{B},\sigma^{2}\}$. This integral has no closed form in general because the Gaussian prior on $\bm{\lambda}_{j}$ is not conjugate to the multinomial likelihood, so we treat $\{\bm{\lambda}_{j}\}_{j=1}^{J}$ as missing data and then estimate $\bm{\theta}=\{\bm{\mu}, \bm{\Phi}, \mathbf{B},\sigma^{2}\}$ by EM algorithm. The standard approach in generalized linear mixed models is the Laplace approximation (\citep{Breslow1993}, \citep{Tierney1986}), which replaced the integral with a Gaussian approximation at the posterior mode. To provide a more readable structure for the following sections, we introduce the basic notations and equalities for standard Laplace approximation.

Taking the logarithm of likelihood~\ref{eq:log-lik-chap2} for each group $j$, we denote $h\left(\bm{\lambda}\right)$ as the logarithm of the integrand. The Laplace  method approximates the integral by forming a second-order Taylor expansion of $h\left(\bm{\lambda}\right)$ around its mode $\hat{\bm{\lambda}} = \arg\max_{\bm{\lambda}}\ell(\bm{\lambda})$. Let $\mathbf{H}$ be the negative Hessian matrix evaluated in the mode and $\mathbf{g}\left(\bm{\lambda}\right)$ be the gradient, then, 
\begin{equation}
    \mathbf{H}(\bm{\lambda}):=-\nabla^{2}h(\bm{\lambda})=\sum_{i\in\mathcal{I}_{j}}y_{i\cdot}\left(\text{diag}(\bm{p}_{i})-\bm{p}_{i}\bm{p}_{i}^{\top}\right)+\bm{\Sigma}^{-1},
    \label{eq: hessian-equality}
\end{equation}
\begin{equation}
    \mathbf{g}\left(\bm{\lambda}\right):=-\nabla h\left(\bm{\lambda}\right)=\sum_{i\in \mathcal{I}_{j}}\bigl(\bm{y}_{i}-y_{i\cdot}\bm{p}_{i}(\bm{\lambda})\bigr)-\bm{\Sigma}^{-1}\bm{\lambda}.
    \label{eq:gradient}
\end{equation}
For the simplicity, we provide three definitions in advance. Define
\begin{equation*}
    d_{q}\left(\bm{\lambda}\right):=\sum_{i\in\mathcal{I}_{j}}y_{i\cdot}p_{iq}\left(\bm{\lambda}\right), \qquad \tilde{\mathbf{D}}\left(\bm{\lambda}\right):=\operatorname{diag}(\mathbf{d}\left(\bm{\lambda}\right)+\sigma^{-2}\mathbf{I}_{\text{Q}}),
\end{equation*}
and the weight matrix $\mathbf{P}\left(\bm{\lambda}\right)\in \mathbb{R}^{Q\times I_{j}}$ with the $i$-th row $\sqrt{y_{i\cdot}}\bm{p}^{\top}_{i}\left(\bm{\lambda}\right)$, then the dense product can be rewritten as,
\begin{equation}
    \mathbf{P}^{\top}\mathbf{P}=\sum_{i\in\mathcal{I}_{j}}y_{i\cdot}\bm{p}_{i}\bm{p}_{i}^{\top},
    \label{eq:weight-prob}
\end{equation}
and finally define $\mathbf{M}_{K}:=\mathbf{I}_{K}+\sigma^{-2}\mathbf{B}^{\top}\mathbf{B}$.

\section{Woodbury Formula}
\label{sec:woodbury}

We introduce the Woodbury matrix identity, which we use throughout this section to invert matrices with a low-rank-plus-diagonal structure, such as $\bm{\Sigma}$ and the exact Hessian $\mathbf{H}$. For any invertible matrix $\mathbf{A}$ and matrices $\mathbf{U}, \mathbf{V}$ of compatible size, the identity reads
\begin{equation}
    \left(\mathbf{A} + \mathbf{U}\mathbf{V}^{\top}\right)^{-1} = \mathbf{A}^{-1} - \mathbf{A}^{-1}\mathbf{U}\left(\mathbf{I} + \mathbf{V}^{\top}\mathbf{A}^{-1}\mathbf{U}\right)^{-1}\mathbf{V}^{\top}\mathbf{A}^{-1}.
    \label{eq:standard-woodbury}
\end{equation}
In our setting the correction has rank $I_j+K$ rather than rank $K$ alone, since both the likelihood and the prior contribute a low-rank piece, so the Woodbury formula reduces the cost of inverting the Hessian from $\mathcal{O}(Q^3)$ to $\mathcal{O}\big(Q(I_j+K)^2+(I_j+K)^3\big)$, which is linear in $Q$ whenever the group size $I_j$ and the factor rank $K$ stay small relative to $Q$. By~\eqref{eq: hessian-equality} and \eqref{eq:weight-prob}, the Hessian can be written as the sum of a likelihood part and a prior part, where the likelihood part is
\begin{equation}
    \sum_{i\in\mathcal{I}_{j}}y_{i\cdot}\left(\text{diag}(\bm{p}_{i})-\bm{p}_{i}\bm{p}_{i}^{\top}\right)=\operatorname{diag}\left(\mathbf{d}\right)-\mathbf{P}^{\top}\mathbf{P},
    \label{eq:likelihood-part}
\end{equation}
with $d_q:=\sum_{i\in\mathcal{I}_j}y_{i\cdot}p_{iq}$. The prior part, transformed by the Woodbury formula, is
\begin{equation}
    \bm{\Sigma}^{-1} = \left(\mathbf{B}\mathbf{B}^{\top} + \sigma^2 \mathbf{I}_Q\right)^{-1} = \sigma^{-2}\mathbf{I}_Q - \sigma^{-4}\mathbf{B}\mathbf{M}_K^{-1}\mathbf{B}^{\top}.
    \label{eq:prior-part}
\end{equation}
Adding \eqref{eq:likelihood-part} and \eqref{eq:prior-part} together, we combine the diagonal terms into $\tilde{\mathbf{D}}:=\operatorname{diag}\left(\mathbf{d}\right)+\sigma^{-2}\mathbf{I}_{Q}$, while the two low-rank terms combine into
\begin{equation}
    \mathbf{P}^{\top}\mathbf{P} + \sigma^{-4}\mathbf{B}\mathbf{M}_K^{-1}\mathbf{B}^{\top}
    = \begin{bmatrix} \mathbf{P}^{\top} & \mathbf{B} \end{bmatrix}
    \begin{pmatrix} \mathbf{I}_{I_{j}} & \bm{0}\\ \bm{0} & \sigma^{-4}\mathbf{M}_K^{-1} \end{pmatrix}
    \begin{pmatrix} \mathbf{P} \\ \mathbf{B}^{\top} \end{pmatrix},
    \label{eq:combined}
\end{equation}
i.e., $\mathbf{U}\mathbf{C}\mathbf{U}^{\top}$. This gives \eqref{eq:hessian-decomp}. Positive semi-definiteness follows from $\mathbf{C}\succ 0$, and the rank does not exceed $I_{j}+K$, the number of columns of $\mathbf{U}$. Two points are worth noting here. First, the cross term $-\bm{p}_{i}\bm{p}_{i}^{\top}$ in the likelihood is exactly what makes $\mathbf{H}$ dense. Its row rank is only 1, and combined over $I_{j}$ observations the rank is at most $I_{j}$. Second, the low-rank part of the prior originates from the factor structure itself, with rank $K$. The two happen to fold into the same $\mathbf{U}$, and this is the entire source of the speed-up in this section. The "denseness" of the data and the "low-rankness" of the prior are, here, the same thing.

Thus, for any $\bm{\lambda}\in\mathbb{R}^{Q}$, define
\begin{equation}
    \mathbf{U}\left(\bm{\lambda}\right):=\left[\mathbf{P}\left(\bm{\lambda}\right)^{\top} \quad \mathbf{B}\right]\in \mathbb{R}^{Q\times (I_{j}+K)}, \qquad \mathbf{C}:=\begin{pmatrix}
    \mathbf{I}_{I_{j}}
    &
    \bm{0}
    \\
    \bm{0}
    &
    \sigma^{-4}\mathbf{M}_{K}^{-1}
    \end{pmatrix}.
    \label{eq:wb-decomp}
\end{equation}
Then the exact Hessian structure satisfies
\begin{equation}
    \mathbf{H}=\tilde{\mathbf{D}}-\mathbf{U}\mathbf{C}\mathbf{U}^{\top}.
    \label{eq:hessian-decomp}
\end{equation}
Applying \eqref{eq:standard-woodbury} to $\mathbf H=\tilde{\mathbf D}+\mathbf U(-\mathbf C)\mathbf U^\top$ gives the Woodbury inverse
\begin{equation}
    \mathbf{H}^{-1}=\left(\tilde{\mathbf{D}}-\mathbf{U}\mathbf{C}\mathbf{U}^{\top}\right)^{-1} =\tilde{\mathbf{D}}^{-1} + \tilde{\mathbf{D}}^{-1}\mathbf{U}\left(\mathbf{C}^{-1} - \mathbf{U}^{\top}\tilde{\mathbf{D}}^{-1}\mathbf{U}\right)^{-1}\mathbf{U}^{\top}\tilde{\mathbf{D}}^{-1}.
    \label{eq:hessian-decomp-wb}
\end{equation}

Let
\begin{equation}
    \Delta\left(\bm{\lambda}\right):= \mathbf{C}^{-1} - \mathbf{U}^{\top}\tilde{\mathbf{D}}^{-1}\mathbf{U}
= \begin{pmatrix}
\mathbf{I}_{I_{j}} - \mathbf{P}\tilde{\mathbf{D}}^{-1}\mathbf{P}^{\top} & -\mathbf{P}\tilde{\mathbf{D}}^{-1}\mathbf{B} \\
-\mathbf{B}^{\top}\tilde{\mathbf{D}}^{-1}\mathbf{P}^{\top} & \sigma^4\mathbf{M}_K - \mathbf{B}^{\top}\tilde{\mathbf{D}}^{-1}\mathbf{B}
\end{pmatrix}
    \label{eq:capacitance}
\end{equation}
denote the capacitance matrix, of the smaller size $(I_{j}+K)\times (I_{j}+K)$ rather than the Hessian's $Q\times Q$, so that $\mathbf{H}^{-1}=\tilde{\mathbf{D}}^{-1}+\tilde{\mathbf{D}}^{-1}\mathbf{U}\Delta^{-1}\mathbf{U}^{\top}\tilde{\mathbf{D}}^{-1}$. For the inverse of $\mathbf H$ to exist and for its Cholesky factorization to be valid, $\Delta$ must be positive definite. This is not an extra condition to check. $\tilde{\mathbf D}\succ 0$ trivially, and the log-posterior is strictly concave at the mode, so $\mathbf H\succ 0$. The matrix determinant lemma applied to \eqref{eq:hessian-decomp-wb} then forces $\Delta\succ 0$ as well.

Putting the pieces together, computing $\mathbf{H}^{-1}$, or a solve $\mathbf{H}^{-1}\bm v$, never requires factorizing a $Q\times Q$ matrix. We only need $\tilde{\mathbf D}^{-1}$, which is diagonal, the $(I_j+K)\times(I_j+K)$ matrix $\Delta(\bm\lambda)$, and a Cholesky factorization of $\Delta$. Since $I_j$ and $K$ are both small and fixed while $Q$ grows, this gives the $\mathcal O\big(Q(I_j+K)^2+(I_j+K)^3\big)$ cost stated at the start of this section, in place of the dense $\mathcal O(Q^3)$ inversion.

\section{Poisson-Multinomial Equivalence}
\label{sec:PME}
    
The Poisson-Multinomial Equivalence established by \citet{Baker1994} states that the Multinomial likelihood conditional on $y_{ij\cdot}$ is exactly equivalent to a product of independent Poisson likelihoods sharing a common additive offset,

\begin{equation}
    \xi_{ij} = \log y_{ij\cdot} - \log\sum_{q=1}^Q \exp(\eta_{ijq}),
    \label{eq:xi_baker}
\end{equation}

The joint Multinomial draw $\bm{y}_{ij} \mid \bm{\lambda}_{j} \sim \mathrm{Multinomial}(y_{ij\cdot}, \bm{p}_{ij})$ is equivalent to $y_{ijq} \mid \bm{\lambda}_{j} \sim \mathrm{Poisson}(\exp(\eta_{ijq} + \xi_{ij}))$ independently for $q = 1, \ldots, Q$, conditional on $y_{ij\cdot}$. We give the definition of $\eta_{ijq}$ and $\bm{p}_{ij}$. The offset~\eqref{eq:xi_baker} absorbs the softmax normalizer and decouples the $Q$ Poisson problems, thus we can solve these $Q$ Poisson independently. 

\citet{Lee2017} extended the framework of \citet{Baker1994} by treating the aggregate count $y_{ij\cdot}$ as Poisson-distributed given the group effect $\bm{\lambda}_{j}$, rather than fixing it as the conditioning variable in the Multinomial likelihood. Marginalizing over this random aggregate recovers the same factorization into independent Poisson components across categories $q = 1, \dots, Q$. As a result, $\xi_{ij}$ becomes the rate parameter of this induced Poisson distribution instead of a fixed conditioning quantity.

In our Multinomial dataset and EM algorithm, Q categories can be decomposed into Q paralleled GLM problems in M-step.

\section{Edgeworth Correction for Group Effects}
\label{subsec:corrected-lambda}

The Laplace E-step needs the posterior mean of $\bm{\lambda}_j$ for each group $j$. But it only computes the posterior mode $\hat{\bm{\lambda}}_j$ and the curvature $\hat{\mathbf{S}}_j$ at that mode. Using the mode instead of the mean is exact when the posterior is Gaussian, since the two coincide. Our posterior is not Gaussian. The multinomial likelihood is skewed in $\bm{\lambda}$, so the mean sits away from the mode. The mode alone therefore feeds a shifted location into every M-step. \citet{katsevich2024} correct the same problem for general Bayesian posteriors, by adding a skew term to the Gaussian factor of the Laplace approximation. We adapt their idea to the multinomial case.

For group $j$, write $\mathrm{lse}(\bm{\eta}_i) := \log\sum_{q=1}^Q \exp(\eta_{iq})$. Collect the third derivatives of the log-likelihood at the mode $\hat{\bm\lambda}_j$ in the tensor
\begin{equation}
    \mathbf{T}_{j,qrs} = \sum_{i \in I_j} y_{ij\cdot} \frac{\partial^3\,\mathrm{lse}(\bm{\eta}_i)}{\partial\eta_q\,\partial\eta_r\,\partial\eta_s}.
    \label{eq:third-tensor}
\end{equation}
The prior is Gaussian, so this tensor captures all of the posterior's skewness. Following \citet{katsevich2024}, the leading-order correction to the mode is
\begin{equation}
    \bm{\mu}_{1,j} = -\tfrac{1}{2}\,\hat{\mathbf{S}}_j\big\langle \mathbf{T}_j, \hat{\mathbf{S}}_j \big\rangle, \qquad \big\langle \mathbf{T}_j, \mathbf{S}\big\rangle_q = \sum_{r,s} \mathbf{T}_{j,qrs}\,\mathbf{S}_{rs}.
    \label{eq:edgeworth}
\end{equation}
We use $\hat{\bm\lambda}_j + \bm\mu_{1,j}$ in place of the raw mode wherever the E-step needs the posterior mean. The remaining error is two orders smaller in the usual Edgeworth bookkeeping, because the intermediate-order term vanishes by Gaussian parity.

For the multinomial likelihood the contraction has a closed form. It avoids ever forming the $Q\times Q\times Q$ tensor:
\begin{equation}
    \big\langle \mathbf{T}_j, \mathbf{S}\big\rangle_q = \sum_{i \in I_j} y_{ij\cdot}\, p_{iq}\Big[\mathbf{S}_{qq} - 2(\mathbf{S}\bm{p}_i)_q - \bm{p}_i^\top\mathrm{diag}(\mathbf{S}) + 2\,\bm{p}_i^\top\mathbf{S}\bm{p}_i\Big],
    \label{eq:T-contraction-multinomial}
\end{equation}
where $\bm p_i$ is observation $i$'s fitted probability vector at the mode. The cost per observation is one matrix-vector product $\mathbf{S}\bm p_i$, so the correction is cheap. We checked \eqref{eq:third-tensor} against finite differences, and \eqref{eq:T-contraction-multinomial} against a brute-force tensor contraction. Both agree to numerical precision.

Intuitively, $\mathbf{T}$ measures how skewed the log-likelihood is at the mode. $\hat{\mathbf S}$ measures how much posterior spread there is to feel that skewness with. Equation~\eqref{eq:edgeworth} converts the two into a shift of the mode toward the heavier tail. The correction is largest for groups with small counts or very unbalanced $\bm p_i$, exactly where the multinomial posterior is most skewed.

\iffalse
In practice this correction is a modest fix, not a decisive one. In oracle single-round tests it recovers only part of the per-round bias in $\sigma^2$. Used on its own, it can also destabilize the EM iterations. Feeding $\hat{\bm\lambda}_j + \bm\mu_{1,j}$ directly into the M-step reintroduces the non-monotonicity problem from the previous section. The algorithm can then fail to reach the convergence tolerance, even though the underlying estimate has improved. Motivated by \citet{Yue2026}, we
\fi 

\section{Laplace-EM Algorithm via Woodbury}
\label{sec: EM-wb}

In this section, we construct a Laplace EM algorithm to estimate the parameters in the FA-DMR. The E-step relies on a Laplace approximation to compute the approximate models of the random effects. \citet{Rubin1982} established the Rubin-Thayer algorithm for the extraction of $\mathbf{B}$ and $\mathbf{D}$ from $\bm{\Sigma}$. Under our model assumption, the idiosyncratic matrix determined by $\sigma^{2}$ rather than a diagonal matrix with different entries. Thus, we introduce one classical result in probabilistic principle component analysis (PPCA) by \citet{Tipping1999}, who had shown that,
\begin{equation}
    \mathbf{B}=\mathbf{U}_{K}\left(\bm{\Lambda}_{K}-\sigma^{2}\mathbf{I}\right)^{1/2}\mathbf{R}.
    \label{eq:tipping-ppca_b}
\end{equation}
For the estimated $\mathbf{B}$, $\sigma^{2}$ satisfies,
\begin{equation}
    \sigma^{2}=\frac{1}{Q-K}\sum_{n=K+1}^{Q}\lambda_{n},
    \label{eq:tipping-ppca_sigma}
\end{equation}
where the $K$ column vectors in the $Q\times K$ matrix $\mathbf{U}_{K}$ are the principal eigenvectors of the sample covariance matrix of the group index, with corresponding eigenvalues $\lambda_{1},\cdots,\lambda_{K}$ in the $K\times K$ diagonal matrix $\bm{\Lambda}_{K}$, and $\mathbf{R}$ is an arbitrary $K\times K$ orthogonal rotation matrix.

The M-step updates model parameters via Rubin-Thayer algorithm \citep{Rubin1982}. For the high dimensional data with large group number $J$ and category number $Q$, we introduce the Woodbury equality to reduce the calculation complexity, and increasingly reduce the runtime rather than multinomial itself. Since Woodbury is the exact surrogate for the matrices from Multinomial distribution, the estimation accuracy should be the same with Multinomial without optimization improvement, but reduce the runtime shortly. 

\paragraph*{E-step: Estimate the Random Effect via LA.}
For each group $j$, the log-posterior of $\bm{\lambda}_{j}$ is given by:
\begin{equation}
    \ell(\bm{\lambda}_{j})=\sum_{i\in \mathcal{I}_{j}}\Bigl[\sum_{q}y_{ijq}\eta_{ijq}(\bm{\lambda}_{j})-y_{ij\cdot}\log \sum_{q}\exp\bigl(\eta_{ijq}(\bm{\lambda}_{j})\bigr)\Bigr]-\frac{1}{2}\bm{\lambda}_{j}^{\top}\bm{\Sigma}^{-1}\bm{\lambda}_{j},
    \label{eq: estep-loglik}
\end{equation}
up to an additive constant independent of $\bm{\lambda}_{j}$, where 
\begin{equation}
    \eta_{ijq}(\bm{\lambda}_{j})=\mu_{q}+\bm{v}_{ij}^{\top}\bm{\phi}_{q}+\lambda_{jq},
    \label{eq: individualcovariate}
\end{equation}
is evaluated separately for each $i\in \mathcal{I}_{j}$, gives the individual-specific covariate structure.
The gradient and Hessian of \eqref{eq: estep-loglik} are~\eqref{eq: hessian-equality} and \eqref{eq:gradient}. Since $\mathbf{P}_{ij}\succeq 0$ and $\bm{\Sigma}^{-1}\succ 0$, we have $\mathbf{H}(\bm{\lambda}_{j})\prec 0$ everywhere, so $\ell(\bm{\lambda}_{j})$ is strictly concave and has a unique maximizer,
\begin{equation}
     \hat{\bm{\lambda}}_{j}
     = \arg\max_{\bm{\lambda}\in \mathbb{R}^{Q}}
     \ell_{j}(\bm{\lambda}_{j}),
     \label{eq:maximizer}
\end{equation}
which can be found by Newton's method. Strict concavity guarantees convergence to $\hat{\bm{\lambda}}_{j}$ from any starting point. The Laplace approximation then replaces the intractable posterior with the Gaussian matched to the model and curvature at,
\begin{equation}
    \bm{\lambda}_{j}^{(t+1)}=
    \bm{\lambda}_{j}^{(t)}
    -s^{(t)}\mathbf{H}\bigl(\bm{\lambda}_{j}^{(t)}
    \bigr)^{-1}\nabla_{\bm{\lambda}_{j}}
    \ell \bigl(\bm{\lambda}_{j}^{(t)}\bigr),
    \label{eq:newton-method}
\end{equation}
The Laplace approximation then yields,
\begin{equation}
    p(\bm{\lambda}_{j}\mid \bm{y}_{ij}) \approx \mathcal{N}\bigl(\hat{\bm{\lambda}}_{j}, \hat{\mathbf{S}}_{j}\bigr), \quad \hat{\mathbf{S}}_{j}=\bigl(-\mathbf{H}(\hat{\bm{\lambda}_{j}})\bigr)^{-1}.
    \label{eq:surrogate-eq}
\end{equation}
In E-step, we also designed the improvement using Multinomial-Poisson Equivalence in the calculation of Hessian before, where the first term with complexity $\mathcal{O}\left(Q^{3}\right)$ in~\ref{eq: hessian-equality} was been removed. However, we observed severe instability on the value of log-likelihood, leading challenging issue of non-monotonicity and optimization. Meanwhile, the Heywood problem arises. We may encounter a severe Heywood problem, where the estimated error covariance matrix collapses during the iterations, with the final estimates degenerating toward zero. Therefore, we use Newton's method with Armijo backtracking search with step size $s^{(t)}\in (0,1]$ for the estimation of $\bm{\lambda}$. 

\paragraph*{M-step: Maximization and Estimation of Fixed Effects.}
With $\hat{\bm{\lambda}}_{j}$ and $\hat{\mathbf{S}}_{j}$ estimated, we solve the maximization problem independently, 
\begin{equation}
    (\hat{\mu}_{q}, \hat{\bm{\phi}}_{q})=\arg \max_{\mu_{q}, \bm{\phi}_{q}}\sum_{i}\bigl[y_{ijq}(\mu_{q}+\bm{v}_{ij}^{\top}\bm{\phi}_{q})-\exp (\mu_{q}+\bm{v}_{ij}^{\top}\bm{\phi}_{q}+\hat{\lambda}_{jq}+\xi_{ij})\bigr].
\end{equation}
This is a Poisson log-linear regression with offset $\lambda_{jq}+\xi_{ij}$, solved by iteratively reweighted least squares. The Q problems are independent and may be parallelized. The EM sufficient statistic for $\bm{\Sigma}$ follows directly from the Laplace output:
\begin{equation}
    \hat{\bm{\Sigma}}_{\text{obs}}=\frac{1}{J}\sum_{j}\mathbb{E}[\bm{\lambda}_{j} \bm{\lambda}_{j}^{\top}\mid \bm{y}_{ij}]\approx \frac{1}{J}\sum_{j}(\hat{\bm{\lambda}}_{j} \hat{\bm{\lambda}}_{j}^{\top}+\hat{\mathbf{S}}_{j}).
    \label{eq:sigma-obs}
\end{equation}
Critically, $\hat{\bm{\Sigma}}_{\text{obs}}$ is fully determined by current Laplace output $(\hat{\bm{\lambda}}_{j}, \hat{\mathbf{S}}_{j})$ and does not need to be re-evaluated during the Rubin-Thayer iterations, thus eliminating the circular dependency inherent in factor-space formulations.
Given $\hat{\bm{\Sigma}}_{\text{obs}}$, we decompose $\bm{\Sigma}=\mathbf{BB}^{\top}+\sigma^2\mathbf{I}_{Q}$ via~\ref{eq:tipping-ppca_b}-\ref{eq:tipping-ppca_sigma} \citep{Tipping1999}. The PLT constraint is enforced post-hoc via a QR decomposition of the top K rows of $\hat{\mathbf{B}}_{\text{new}}$, followed by sign correction to ensure positive diagonal entries and zeroing of the upper triangle. Applying PLT only after the Rubin–Thayer sub-iterations converge, rather than at each inner step, avoids the discontinuities that would otherwise disrupt convergence.

Given $\hat{\mathbf{B}}$ and $\hat{\sigma}^2$, the factor scores are recovered by the Bayes linear estimator \citep{Gelman2013}:
\begin{equation}
\hat{\mathbf{f}}_j = \left( \mathbf{I}_K + \frac{1}{\hat{\sigma}^2} \hat{\mathbf{B}}^\top \hat{\mathbf{B}} \right)^{-1} \frac{1}{\hat{\sigma}^2} \hat{\mathbf{B}}^\top \hat{\boldsymbol{\lambda}}_j ,
\label{eq: Bayeslinearestimator}
\end{equation}
requiring only $K\times K$ inversion.

% Chapter~\ref{cha:methodology} derived a correct and convergent Laplace-EM algorithm, but left its dominant per-iteration cost unaddressed. Each Newton step in E-step requires forming and solving the $Q\times Q$ system~\eqref{eq: hessian-equality} at cost $\mathcal{O}(I_{j}Q^{2}+Q^{3})$, in the meanwhile, each inner iteration of the M-step factor update~\eqref{eq: Rubin3}-\eqref{eq: Rubin4} requires forming $\bm{\Sigma}^{-1}$ densely at cost $\mathcal{O}(Q^{3})$. For large Q, these costs are not merely inconvenient but make the algorithm impractical. This chapter develops an exact acceleration method to remove this bottleneck by showing that the Poisson-Multinomial equivalence that already underlies the M-step can be redeployed via auxiliary-variable construction, which can remove the cross-category coupling from the E-step Hessian as well.  

We give a runtime comparison table for a small-scale test with moderate $J$ and $Q$, however, the standard multinomial and \texttt{glmmTMB} packages cost more than expected compared with EM-algorithm via Woodbury formula, which are shown below,

\begin{table}[ht]
\centering
\begin{tabular}{llccccc}
\toprule
Setting & Method & Time (s) & Speedup & $\hat\sigma^2$ & $d(\hat B, B_0)$ & Converged \\
\midrule
\multirow{4}{*}{$J=50,\ Q=50$}
 & std             & 55.58  & 1.00$\times$ & 0.2895 & 0.3266 & NO \\
 & Woodbury   & 24.56  & 2.26$\times$ & 0.2895 & 0.3266 & NO \\
 & Wb + Corr       & 25.47  & 2.18$\times$ & 0.2997 & 0.3260 & NO \\
 & glmmTMB    & 40.10  & 1.39$\times$ & 0.3874 & 0.2536 & yes \\
\midrule
\multirow{4}{*}{$J=50,\ Q=100$}
 & std             & 83.90  & 1.00$\times$ & 0.2739 & 0.6310 & NO \\
 & Woodbury   & 37.03  & 2.27$\times$ & 0.2739 & 0.6310 & NO \\
 & Wb + Corr       & 38.89  & 2.16$\times$ & 0.2877 & 0.6318 & NO \\
 & glmmTMB    & 130.29 & 0.64$\times$ & 0.4056 & 0.3145 & yes \\
\midrule
\multirow{4}{*}{$J=100,\ Q=50$}
 & std        & 114.44 & 1.00$\times$ & 0.2935 & 0.5317 & NO \\
 & Woodbury   & 48.77  & 2.35$\times$ & 0.2935 & 0.5317 & NO \\
 & Wb + Corr  & 50.27  & 2.28$\times$ & 0.3033 & 0.5308 & NO \\
 & glmmTMB    & 108.95 & 1.05$\times$ & 0.4201 & 0.4306 & yes \\
\midrule
\multirow{4}{*}{$J=100,\ Q=100$}
 & std        & 167.49 & 1.00$\times$ & 0.2768 & 0.2899 & NO \\
 & Woodbury   & 71.49  & 2.34$\times$ & 0.2768 & 0.2899 & NO \\
 & Wb + Corr  & 72.68  & 2.30$\times$ & 0.2914 & 0.2898 & NO \\
 & glmmTMB    & 412.37 & 0.41$\times$ & 0.3920 & 0.2193 & yes \\
\bottomrule
\end{tabular}
\caption{Runtime and estimation accuracy across four $(J,Q)$ settings, seed=1, true $\sigma^2=0.30$. Speedup is relative to \texttt{std}. \texttt{std} and \texttt{Woodbury} agree exactly on $\hat\sigma^2$ and $d(\hat B, B_0)$ by construction (same algebra); the gap in runtime isolates the Woodbury acceleration itself.}
\label{tab:speed-accuracy}
\end{table}

Table~\ref{tab:speed-accuracy} reports runtime and estimation accuracy across four $(J,Q)$ settings. This speedup is a pure acceleration effect. The Woodbury and dense estimators agree exactly on $\hat\sigma^2$ and $d(\hat B,B_0)$ in all four settings, since the two implementations solve the same optimization problem through different linear algebra.

The dense, Woodbury, and corrected estimators did not converge within the fixed iteration budget in any of the four settings. This is expected, not a failure. The approximate E-step voids the classical EM monotonicity guarantee, so the log-evidence trajectory is not monotone. A fixed-budget run under a non-monotone objective need not satisfy a convergence tolerance, and does not indicate that the estimate is unreliable. \texttt{glmmTMB} converges in every setting, but its runtime advantage does not hold uniformly: it is faster than the dense baseline at $Q=50$ but substantially slower at $Q=100$. More importantly, glmmTMB attains a lower $d(\hat{\mathbf{B}}, \mathbf{B}_0)$ than our uncorrected estimator in all four settings. This gap in $B$-recovery, not runtime, is the primary shortcoming of the base algorithm and motivates the correction developed in the next section.

\subsection{Limitations of the Standard EM-Algorithm \& Solutions}

Although the runtime of Woodbury-EM algorithm reduces increasingly compared to standard EM algorithm and existing packages in R, i.e., \texttt{glmmTMB}, the accuracy is challenged, because our previous work all based on the acceleration rather than optimization. There are some facts that interesting when we did the small-scale tests, that the estimated error of loading matrix $\mathbf{B}$ is accumulated through iterations due to the nature of EM algorithm, resulting the final estimation of $\mathbf{B}$ in slightly less accuracy than $\texttt{glmmTMB}$ and $\texttt{gllvm}$. Meanwhile, the estimation of $\mathbf{B}$ affects $\sigma^{2}$ in this iteration, according to~\ref{eq:tipping-ppca_b}--\ref{eq:tipping-ppca_sigma}, we refer to this as coupling of two steps. 

For more details, we come back to the prior in E-step $\bm{\Sigma}^{-1}=\left(\hat{\mathbf{B}}\hat{\mathbf{B}}^{\top}+\sigma^{2}\mathbf{I}\right)^{-1}$, which affected by $\hat{\mathbf{B}}^{(t)}$ itself in $t$-th iteration. In practice, this point leads to a self-confirming feedback loop. Along the directions that have already been identified as important by $\hat{\mathbf{B}}$, the prior variance is large and there is less shrinkage, while in the orthogonal directions, the prior variance is only $\sigma^{2}$ and there is more shrinkage. Therefore, $\hat{\bm{\lambda}}_{j}$ mainly spreads along the directions that have already been identified by $\hat{\mathbf{B}}$, and the M-step then uses these $\hat{\lambda}_j$ to estimate $\hat{\mathbf{B}}$ back to the same directions. The systematic localization experiments, including two comparison groups with $\bm{\mu},\bm{\phi}$ fixed and $\mathbf{B},\sigma^2$ varying, and with $\mathbf{B},\sigma^2$ fixed and $\bm{\mu},\bm{\phi}$ varying, precisely locate the source of the problem in this mechanism, rather than the estimation of $\bm{\mu},\bm{\phi}$ or the Rubin-Thayer method itself.

In the following section, we arrange three methods of correction and compare the results of both accuracy, runtime and robustness.

% Add the result table

\subsubsection{ISO: Isotropic Prior Setting for E-step}
\label{subsubsec:isotropic-prior}

For the self-confirming issue in the estimation of $\mathbf{B}$, we rewrite the prior in E-step using isotropic working prior $\mathcal{N}\left(\bm{0},\tau^{2}\mathbf{I}\right)$ instead of $\mathcal{N}\left(\mathbf{0},\bm{\Sigma}\right)$, where $\tau^{2}:=\operatorname{tr}\left(\hat{\mathbf{B}}\hat{\mathbf{B}}^{\top}\right)/Q+\hat{\sigma}^{2}$. Replacing $\bm{\Sigma}$ by $\tau^{2}\mathbf{I}$ in Newton step and define the surrogate $\hat{\mathbf{S}}_{j}=\left(\mathbf{F}_{j}+\tau^{-2}\mathbf{I}\right)^{-1}$ with Fisher information $\mathbf{F}_{j}=\sum_{i\in\mathcal{I}_{j}}y_{i\cdot}\left(\text{diag}(\bm{p}_{i})-\bm{p}_{i}\bm{p}_{i}^{\top}\right)$ for the $j$-th group. If we assume that $\mathbf{F}_{j}$ can be approximated as $I_{j}\mathbf{I}$, the posterior mode $\hat{\bm{\lambda}}_{j}$ degenerates to scalar shrinkage of the pure-data MLE $\tilde{\bm{\lambda}}_{j}$,
\begin{equation}
    \hat{\bm{\lambda}}_{j}=c\tilde{\bm{\lambda}}_{j}, \qquad c=\frac{I_{j}\tau^{2}}{I_{j}\tau^{2}+1}\in\left(0,1\right).
    \label{eq:MLE-lam}
\end{equation}
Correspondingly, the second moment of M-step is,
\begin{equation}
    \hat{\bm{\Sigma}}^{\text{iso}}\approx c^{2}\bm{\Sigma}+\kappa\mathbf{I},
    \label{eq:iso-M}
\end{equation}
where $\kappa=\frac{c(c+1)}{I_{j}}$. $c^{2}\bm{\Sigma} + \kappa \mathbf{I}$ has exactly the same eigenvectors as $\bm{\Sigma}$, so the subspace is preserved exactly. The cost is that the loadings are shrunk by a factor of $c$ overall, while $\hat{\bm{\sigma}}^2$ is systematically overestimated by $\kappa$. The Heywood problem depends exactly on the condition that the prior contains $\hat{\mathbf{B}}$. In Equation~\eqref{eq:iso-M}, the prior only contains $\tau^2$ and no longer contains $\hat{\mathbf{B}}$, so the first-level feedback loop is physically cut off at this step.

\subsubsection{Debias: Moment-Matching Path}
\label{subsubsec:debias}

The cost of Equation~\ref{eq:iso-M} is not free. The overall shrinkage of the loading matrix by a factor of $c$ and the systematic overestimation of $\hat{\sigma}^2$ by $\kappa$ are both systematic biases that can be calculated exactly. Therefore, these biases should be corrected exactly without introducing another approximation.

Fay-Herriot model~\citep{Herriot1979} introduced a moment-matching estimator, which we adapt for our group random effect $\bm{\lambda}$. They use the weighted average of observation and regression estimation as the final estimation. Instead of starting from true group random effect $\bm{\lambda}$, we begin at the pure-data MLE $\tilde{\bm{\lambda}}$ that is estimated directly from the dataset and assume that $\tilde{\bm{\lambda}}\mid\bm{\lambda}\sim\mathcal{N}\left(\bm{\lambda},\mathbf{F}^{-1}\right)$ for each group, where $\bm{\lambda}$ is known and $\mathbf{F}$ is the Fisher information matrix given by the asymptotic theory of MLE. Adapting the isotropic prior setting from Section~\ref{subsubsec:isotropic-prior}, the weighted average of $\hat{\bm{\lambda}}$ takes the form of the classical Bayes estimator in the case that the prior is known. We drop the group index $j$ when discussing a generic group, it is restored whenever a sum over groups appears. 

\begin{equation}
    \delta^{'}=\frac{A^{*}}{A^{*}+D}Y+\frac{D}{A^{*}+D}Y^{*},
    \label{eq:classical-bayes1}
\end{equation}
where $Y$ denotes the direct estimator and $Y^{*}$ denotes the regression-synthetic estimator~\citep{Herriot1979, Pfeffermann2011}. Furthermore, $A^{*}$ and $D$ denote the variance respectively. Using our notation together with the group working prior $\bm{\lambda}\sim \mathcal{N}\left(\bm{0}, \tau^{2}\mathbf{I}\right)$, the expression can be written as,
\begin{equation}
    \hat{\bm{\lambda}}=\frac{\tau^{2}\mathbf{I}}{\tau^{2}\mathbf{I}+\mathbf{F}^{-1}}\tilde{\bm{\lambda}}+\frac{\mathbf{F}^{-1}}{\tau^{2}\mathbf{I}+\mathbf{F}^{-1}}\bm{0}=\frac{\tau^{2}\mathbf{I}}{\tau^{2}\mathbf{I}+\mathbf{F}^{-1}}\tilde{\bm{\lambda}}.
    \label{eq:classical-bayes2}
\end{equation}
This is the standard form of the classical empirical Bayes shrinkage estimator. Given the sampling variance $\mathbf{F}^{-1}$ and the between-group working variance $\tau^{2}\mathbf{I}$, the shrinkage weight is determined by the proportion of data precision to the total precision from both the data and the prior. The standard expression of the BLUP in mixed models is another way to write the same formula [Citation].

In the implementation, $\mathbf{F}$ for each group does not need to be calculated again. It can be directly obtained from $\hat{\mathbf{S}}=\left(\mathbf{F}+\tau^{-2}\mathbf{I}\right)^{-1}$ as $\mathbf{F}=\hat{\mathbf{S}}^{-1}-\tau^{-2}\mathbf{I}$. Equation~\ref{eq:classical-bayes2} can be reversed to obtain the pure-data MLE,
\begin{equation}
\tilde{\bm{\lambda}}=\mathbf{F}^{-1}\hat{\mathbf{S}}^{-1}\hat{\bm{\lambda}}.
\label{eq:pure-data-lambda}
\end{equation}
Since $\bm{\lambda}$ and $\tilde{\bm{\lambda}}$ are two different points, let $\bm{\epsilon}:=\tilde{\bm{\lambda}}-\bm{\lambda}$ denote the sampling noise contained in the pure-data MLE with the variance according to the Cramer-Rao bound $\operatorname{Var}\left(\bm{\epsilon}\right)=\mathbf{F}^{-1}$. Therefore, the between-group covariance can be estimated by
\begin{equation}
\hat{\bm{\Sigma}}=\frac{1}{J}\sum_{j=1}^{J}\left(\tilde{\bm{\lambda}}_j\tilde{\bm{\lambda}}_j^{\top}-\mathbf{F}_j^{-1}\right).
\label{eq:debias-Sigma}
\end{equation}
The estimation in iteration does not contain $\hat{\mathbf{B}}$. For each group, we subtract its own sampling variance $\mathbf{F}_j^{-1}$, which does not depend on other groups or on the current estimate of $\hat{\mathbf{B}}$. This gives a direct way to remove the measurement error variance and recover the variance of the underlying distribution.

The scalar approximation in Equation~\eqref{eq:iso-M} is only used to explain the form of this step under an ideal setting. In the actual implementation, Equation~\eqref{eq:debias-Sigma} uses the group-specific $\mathbf{F}_j$ for each group, which can be anisotropic and does not depend on this approximation. For the multinomial model, $\mathbf{F}_j$ is exactly singular in the all-ones direction, since adding the same constant to all categories does not change the softmax probabilities. Therefore, a small ridge term with $\texttt{ridge}=10^{-3}$ is added to the inversion step in the calculation of $\hat{\mathbf{S}}$ in the implementation.

\paragraph*{Phenomenon.} A problem was been found in both low-information settings ($Q=20,J=30,N_{\mathrm{per\ group}}=8$) and large-scale settings ($Q=200,J=160,N_{\mathrm{per\ group}}=15$). The debiased $\hat{\sigma}^{2}$ directly reaches the numerical lower bound $10^{-8}$ used in the closed-form implementation, regardless of the random seed. The two experiments are shown below,

\begin{table}[H]
\centering
\begin{tabular}{lcccc}
\toprule
Setting ($\sigma^{2}_{\text{true}}=0.3$) & Method & $\hat\sigma^2$ & B distance & Tucker \\
\midrule
\multirow{4}{*}{$Q=20,\ K=2,\ J=30,\ I=8$}
 & std    & 0.2977 & 0.4631 & 0.9031 \\
 & iso    & 0.3867 & 0.2818 & 0.9739 \\
 & debias & $1\times10^{-8}$ & 0.2819 & 0.9749 \\
 & FH-EM   & 0.3187 & 0.2753 & 0.9740 \\
\midrule
\multirow{4}{*}{$Q=200,\ K=2,\ J=160,\ I=15$}
 & std    & 0.3028 & 0.4073 & 0.8784 \\
 & iso    & 0.4461 & 0.1789 & 0.9874 \\
 & debias & $1\times 10^{-8}$ & 0.1701 & 0.9884 \\
 & FH-EM   & 0.3304 & 0.1782 & 0.9872 \\
\bottomrule
\end{tabular}
\caption{Comparison across estimation methods for two problem sizes.}
\label{tab:method-comparison}
\end{table}

\paragraph*{Mechanism.} Equation~\ref{eq:debias-Sigma} subtracts a non-negligible quantity $\mathbf{F}_{j}^{-1}$ from a quantity $\tilde{\bm{\lambda}}_{j}\tilde{\bm{\lambda}}_{j}^{\top}$ that is not necessarily large. When the information from each group is insufficient, this subtraction can easily overshoot, resulting in negative eigenvalues in $\hat{\bm{\Sigma}}$. Once the closed-form PPCA sees that some of the smallest eigenvalues are negative, $\sigma^{2}$ estimation~\eqref{eq:tipping-ppca_sigma} is directly pushed to the numerical lower bound. This is a different mechanism from the first-level feedback loop. In both experiments, the $d_{\mathrm{true}}$ of the debias method is almost identical to that of the isotropic method. This means that the accuracy of the estimated subspace is not affected. The reason is that PPCA separates the estimation of the eigenvectors from the estimation of the smallest eigenvalues. A shift in the magnitude does not rotate the eigenvectors. This is exactly what we define above as the Heywood problem. The collapse occurs in the subtraction step of the moment de convolution itself and does not depend on whether $\hat{\mathbf{B}}$ is included in the prior. The standard method and earlier versions of the algorithm do not show this specific failure mode. It is a new problem introduced by the particular design choice of the debiasing correction.

Honestly, whether this collapse occurs depends on the average count per category, which is approximately proportional to the magnitude of $\mathbf{F}_{j}$. It does not directly depend on $Q$, $J$, or $N{\mathrm{per\ group}}$ themselves. Although $Q$, $J$, and $I_{\mathrm{per\ group}}$ are all much larger in the second experiment than in the first one, the total count range in the simulation, $y_{\mathrm{tot}}\in[200,400]$, was not increased together with $Q$. Therefore, the average count per category was actually more sparse, and the same collapse was reproduced. This reminds us that, when determining whether this failure mode will occur, we should focus on the information per category rather than simply considering the overall data size.

\subsubsection{Fay-Herriot EM Algorithm (FH-EM): Maximum Likelihood Path}

\paragraph*{Model Assumption.} The Fay-Herriot marginal likelihood in Equation~\eqref{eq:fh-marginal-likelihood} borrows its structure directly from the classical small-area model~\citep{Herriot1979, Pfeffermann2011}: a direct estimate observed with a known sampling covariance, linked through a common between-group Gaussian distribution. Transplanting this structure from survey-based small-area estimation to our setting requires several assumptions that do not hold exactly here, and we state them explicitly rather than leave them implicit in Equation~\eqref{eq:observation_model}.

\begin{enumerate}
\item[(A1)] \emph{$\mathbf F_j$ is treated as known, though it is not.} In the classical model, the sampling covariance $\psi_i$ is a design quantity, fixed and known prior to estimation. Here $\mathbf F_j$ is the multinomial Fisher information evaluated at $\hat{\bm\lambda}_j^{\text{iso}}$, the mode of an auxiliary isotropic-prior Laplace approximation (Section~\ref{subsubsec:isotropic-prior}), and is therefore itself an estimated quantity that depends on the current global parameters $(\hat{\mathbf{B}},\hat\sigma^2)$ through $\tau^2$. Equation~\eqref{eq:observation_model} treats $\mathbf F_j$ as fixed and non-stochastic within one EM round; this is the same plug-in approximation already flagged for the marginal likelihood in Equation~\eqref{eq:gaussian-mix}, and it propagates into the E-step and M-step derivations exactly the same way. It is not an additional error introduced by the EM reformulation; it is the same approximation, made explicit at the point where it is used rather than only at the point where its consequence (the Gaussian-mixture correction) is visible.

\item[(A2)] \emph{Validity regime of the Gaussian approximation.} $\tilde{\bm\lambda}_j\mid\bm\lambda_j\sim\mathcal N(\bm 0,\mathbf F_j^{-1})$ relies on the asymptotic normality of the Laplace-approximated MLE. This is accurate when the per-category information within a group is large, and degrades as the average count per category shrinks, independent of $J$ or $Q$ themselves --- the same information-per-category condition identified in the debias failure analysis (Section~\ref{subsubsec:debias}). We do not assume this approximation is exact for any finite sample; we assume it is adequate in the regime our simulations and asymptotic theory (Chapter~\ref{cha:asymptotics}) target.

\item[(A4)] \emph{Identifiability along the singular direction.} $\mathbf F_j$ is exactly singular along the all-ones direction, since a constant shift to all categories leaves the softmax probabilities unchanged. $\mathbf F_j+\bm\Sigma^{-1}$ is well-posed without an additional ridge because $\bm\Sigma^{-1}\succ0$ regularizes this direction on its own; however, $\tilde{\bm\lambda}_j$ itself is only defined up to whatever convention fixes this direction, and that convention is inherited from wherever $\bm\lambda_j$'s own identifiability constraint is imposed, not introduced fresh here.

\item[(A5)] \emph{Scope of the ascent-property claim.} The genuine $\arg\max$ structure and the ascent property established in this section apply to one inner update of $\bm\Sigma$ for a fixed $(\tilde{\bm\lambda}_j,\mathbf F_j)$ pair. Because $\mathbf F_j$ is itself re-estimated every outer round from the evolving isotropic-prior stage, as noted in (A2), the full nested procedure is a generalized EM in the sense that the auxiliary quantities being conditioned on change between rounds. We therefore claim the ascent property only for the inner FH-EM sub-problem, exactly as stated; we do not claim that a single fixed objective function is monotonically increased across outer rounds of the whole isotropic-plus-FH-EM procedure.

\item[(A6)] \emph{What the positive-definiteness guarantee does and does not cover.} $\mathbf S$ in Equation~\eqref{eq:fhem_mstep} is a sum of positive semi-definite terms and therefore cannot produce the negative-eigenvalue collapse that afflicts the debias estimator's moment subtraction (Section~\ref{subsubsec:debias}). This rules out that specific failure mode by construction. It does not imply that $\hat\sigma^2$ can never reach the numerical floor: if the data genuinely favor $\sigma^2\to0$ in the ordinary maximum-likelihood sense, FH-EM has no more immunity to this classical boundary solution than any other exact likelihood-based method, and this is a separate phenomenon from the one FH-EM is designed to correct.

\item[(A7)] \emph{Local, not global, optimum.} As with any EM algorithm applied to a factor-analytic latent-variable model, the likelihood surface need not be unimodal, and the fixed point reached depends on initialization. We do not claim FH-EM recovers the global maximum of Equation~\eqref{eq:fh-marginal-likelihood}; we claim that, given a fixed starting point, each step does not decrease the Q-function in Equation~\eqref{eq:fhem-Q-function}.
\end{enumerate}

Equation~\eqref{eq:debias-Sigma} is a moment matching formula rather than the $\arg\max$ of any objective function. Therefore, the ascent property of the standard EM algorithm does not apply to this step from the beginning. Its validity comes from a separate argument that the moment estimator is consistent, rather than from the EM theory itself. The collapse is a direct consequence of this point. If we want to add a boundary correction to this step to prevent the degeneration, similar to the idea of adjusted ML, there is no natural objective function to modify. It is not maximizing any objective function, so we cannot simply say that a penalty term is added to the objective function.

The first stage of the debias method, which uses the Laplace approximation with an isotropic prior, remains unchanged and still produces $(\tilde{\bm{\lambda}}_j,\mathbf{F}_j)$. In the second stage, we replace the moment matching step with a genuine $\arg\max$ problem. We define $\tilde{\bm{\lambda}}_{j}-\bm{\lambda}_{j}=\bm{\epsilon}_j$ and have the conditional distribution,

\begin{equation} 
    \bm{\epsilon}_{j}\mid \bm{\lambda}_{j} \sim \mathcal{N}\left(\mathbf{0}, \mathbf{F}_j^{-1}\right). 
    \label{eq:observation_model} 
\end{equation} 
Together with $\bm{\lambda}_{j}\sim \mathcal{N}\left(\bm{0},\bm{\Sigma}\right)$, we integrate on $\bm{\lambda}_{j}$ marginally. The density of $\tilde{\bm{\lambda}}_{j}$ is given by $p\left(\bm{\lambda}_{j},\tilde{\bm{\lambda}}_{j}\right)=p\left(\tilde{\bm{\lambda}}_{j}\mid \bm{\lambda}_{j}\right)p\left(\bm{\lambda}_{j}\right)$ and the marginal integral,
\begin{equation}
    p\left(\tilde{\bm{\lambda}}_{j}\right)=\int \mathcal{N}\left( \bm{\lambda}_{j};\,\bm{0},\mathbf{F}_{j}^{-1}\right)\mathcal{N}\left(\bm{\lambda}_{j};\,\bm{0},\bm{\Sigma}\right)\,d\bm{\lambda}_{j},
    \label{eq:gaussian-mix}
\end{equation}
which is an approximation. Strictly speaking, $\mathbf{F}_j$ depends on the realized $\bm{\lambda}_j$ through the multinomial information matrix, so Equation~\eqref{eq:observation_model} is a conditional statement, in which $\mathbf{F}_{j}$ is the function of $\bm{\lambda}_{j}$. Thus, the integral~\eqref{eq:gaussian-mix} is an Gaussian scale mixture, rather than an exact Gaussian. We write it as an approximated formula,
\begin{equation} 
    l_{\mathrm{FH}}(\bm{\Sigma}) \approx \sum_{j=1}^{J} \log \mathcal{N}\left(\tilde{\bm{\lambda}}_j;\, \mathbf{0}, \bm{\Sigma} + \mathbf{F}_j^{-1}\right). 
    \label{eq:fh-marginal-likelihood} 
\end{equation}
The right side is exactly the marginal likelihood of the multivariate Fay-Herriot model~\citep{Herriot1979, Pfeffermann2011}. $\tilde{\bm{\lambda}}_j$ is a direct estimate with known sampling covariance $\mathbf{F}_j^{-1}$, while $\bm{\Sigma}$ is the between-group covariance. When all $\mathbf{F}_j$ are equal, Equation~\eqref{eq:fh-marginal-likelihood} is exactly the MLE of $l_{\text{FH}}(\bm{\Sigma})$. When $\mathbf{F}_j$ are not equal, which naturally arises since the information from different groups can differ and our model allows different $I_j$, the debias estimator in Equation~\eqref{eq:debias-Sigma} remains consistent, and each term $\tilde{\bm{\lambda}}_j\tilde{\bm{\lambda}}_j^{\top}-\mathbf{F}_j^{-1}$ has expectation $\bm{\Sigma}$ regardless of $\mathbf{F}_j$, so the unweighted average converges to $\bm{\Sigma}$ as $J\to\infty$. It is not, however, fully efficient. Maximizing $l_{\text{FH}}(\bm{\Sigma})$ instead weights each group by its precision $(\bm{\Sigma}+\mathbf{F}_j^{-1})^{-1}$, so groups with higher information contribute more to the estimate of $\bm{\Sigma}$. Equal weighting, as used by debias, discards this information and results in a larger asymptotic variance.

\paragraph*{E-step.}
\label{para:fhem-estep}
Equation~\eqref{eq:fh-marginal-likelihood} carries the approximation error discussed above. To avoid it, we introduce a genuine EM algorithm, where the missing data are $\bm{\lambda}_j$ and the observations are $\tilde{\bm{\lambda}}_j$. Given the current $\bm{\Sigma}$ and denote $\mathbf{V}_j = \left(\mathbf{F}_j + \bm{\Sigma}^{-1}\right)^{-1}$ as the covariance matrix of $\bm{\lambda}_{j}\mid \tilde{\bm{\lambda}}_{j}$, the joint log-likelihood is,
\begin{align}
    \log p\left(\bm{\lambda}_{j},\tilde{\bm{\lambda}}_{j}\right)&=-\frac{1}{2}\left(\bm{\lambda}_{j}-\tilde{\bm{\lambda}}_{j}\right)^{\top}\mathbf{F}_{j}\left(\bm{\lambda}_{j}-\tilde{\bm{\lambda}}_{j}\right)-\frac{1}{2}\bm{\lambda}_{j}^{\top}\bm{\Sigma}^{-1}\bm{\lambda}_{j}\\
    &=-\frac{1}{2}\left[\bm{\lambda}_{j}^{\top}\left(\mathbf{F}_{j}+\bm{\Sigma}^{-1}\right)\bm{\lambda}_{j}-2\tilde{\bm{\lambda}}_{j}^{\top}\mathbf{F}_{j}\bm{\lambda}_{j}\right]\\
    &=-\frac{1}{2}\left[\left(\bm{\lambda}_{j} - \mathbf{V}_{j} \mathbf{F}_{j} \tilde{\bm{\lambda}}_{j}\right)^{\top} \mathbf{V}_{j}^{-1} \left(\bm{\lambda}_{j} - \mathbf{V}_{j} \mathbf{F}_{j} \tilde{\bm{\lambda}}_{j}\right) - \tilde{\bm{\lambda}}_{j}^{\top} \mathbf{F}_{j} \mathbf{V}_{j} \mathbf{F}_{j} \tilde{\bm{\lambda}}_{j}\right],
    \label{eq:joint-log-density}
\end{align}
up to a constant independent to $\bm{\lambda}_{j}$. Notice that $\bm{\lambda}_{j}$ only appears in the first term, the posterior therefore is given by $\mathbf{V}_{j}$ and the posterior mean $\mathbf{m}_{j} = \mathbf{V}_{j} \mathbf{F}_{j} \tilde{\bm{\lambda}}_{j}$,
\begin{equation}
    \bm{\lambda}_{j}\mid \tilde{\bm{\lambda}}_{j}\sim \mathcal{N}\left(\mathbf{m}_{j},\mathbf{V}_{j}\right).
    \label{eq:fh-posterior}
\end{equation}
Interestingly, we extend the posterior mean $\mathbf{m}_{j} = \mathbf{V}_{j} \mathbf{F}_{j} \tilde{\bm{\lambda}}_{j}=\left(\mathbf{F}_{j}+\bm{\Sigma}^{-1}\right)^{-1}\mathbf{F}_{j}\tilde{\bm{\lambda}}_{j}$, which is  the similar formula as the shrinkage factor in Equation~\eqref{eq:classical-bayes2} and implies $\mathbf{m}_{j}$ to be the posterior mean of $\bm{\lambda}_{j}$ given $\bm{\Sigma}$.  

We call back the structure of standard EM-algorithm and find the complete log-likelihood. Since $\bm{\Sigma}$ is not contained in $\tilde{\bm{\lambda}}_{j}\mid \bm{\lambda}_{j}$, we have,
\begin{align}
    l_{\text{complete}}\left(\bm{\Sigma}\right)&=\sum_{j}\log \mathcal{N}\left(\bm{\lambda}_{j};\,\bm{0},\bm{\Sigma}\right)=-\frac{J}{2}\log|\bm{\Sigma}| - \frac{1}{2}\sum_{j}\bm{\lambda}_j^{\top}\bm{\Sigma}^{-1}\bm{\lambda}_j + \text{const}\\
    &= -\frac{J}{2}\log|\bm{\Sigma}| - \frac{1}{2}\operatorname{tr}\left(\bm{\Sigma}^{-1}\sum_{j}\bm{\lambda}_j\bm{\lambda}_j^{\top}\right) + \text{const}
    \label{eq:fhem-l-complete}
\end{align}
Applying $\mathbb{E}\left[\mathbf{X}\mathbf{X}^{\top}\right]=\operatorname{Var}\left(\mathbf{X}\right)+\mathbb{E}\left[\mathbf{X}\right]\mathbb{E}\left[\mathbf{X}\right]^{\top}$, we substitute $\mathbb{E}\left[\bm{\lambda}_{j}\bm{\lambda}_{j}^{\top}\mid \tilde{\bm{\lambda}}_{j}\right]=\mathbf{V}_{j}+\mathbf{m}_{j}\mathbf{m}_{j}^{\top}$ into~\eqref{eq:fhem-l-complete} and obtain the Q-function for maximization,
\begin{equation}
    Q(\bm{\Sigma} \mid \bm{\Sigma}^{\text{old}}) = -\frac{J}{2}\log|\bm{\Sigma}| - \frac{1}{2}\operatorname{tr}\left(\bm{\Sigma}^{-1}\sum_{j}\left(\mathbf{m}_j\mathbf{m}_j^{\top} + \mathbf{V}_j\right)\right) + \text{const}.
    \label{eq:fhem-Q-function}
\end{equation}

\paragraph*{M-step.}
\label{para:fhem-mstep}
Calculate the derivative of~\eqref{eq:fhem-Q-function} and obtain the solution,
\begin{equation} \mathbf{S} = \frac{1}{J} \sum_{j=1}^{J} \left( \mathbf{m}_j \mathbf{m}_j^{\top} + \mathbf{V}_j \right). 
\label{eq:fhem_mstep} 
\end{equation}
As the similar as standard Laplace-EM algorithm in Section~\ref{sec: EM-wb}, we apply the closed-form PPCA solution in Equations~\eqref{eq:tipping-ppca_b}-\eqref{eq:tipping-ppca_sigma} to $\mathbf{S}$ to decompose $\mathbf{B}$ and $\sigma^{2}$.

FH-EM algorithm can solve the Heywood problem that happens in debias method~\ref{subsubsec:debias} by establishing a strictly positive definite matrix $\mathbf{S}$. The collapse mechanism is thus structurally blocked rather than simply not appearing in a particular experiment. This is the specific reason why the ``Heywood'' problem is structurally corrected by the new method.

We tested FH-EM algorithm on standard EM-algorithm setting without Woodbury formula usage. Since these optimized processes work well on standard EM algorithm, we can also use the Woodbury to reduce the calculation cost and complexity, which is introduced in the next section. 

\paragraph*{Algorithm.}
\label{para:fhem-algorithm}
Algorithm~\ref{alg:fhem} summarizes the FH-EM procedure derived above. The inputs $\tilde{\bm{\lambda}}_j$ and $\mathbf{F}_j$ come from Stage~1 and remain fixed throughout. Each iteration alternates between the E-step in Equations~\eqref{eq:fh-posterior} and the M-step in Equation~\eqref{eq:fhem_mstep}, followed by a closed-form PPCA update. Both terms in the M-step are positive semi-definite by construction, so $\mathbf{S}$ is always positive definite and the collapse described in Level~2 of Table~\ref{tab:heywood-levels} cannot occur. Convergence is monitored through the marginal log-likelihood $l_{\text{FH}}(\bm{\Sigma})$ rather than the change in $\bm{\Sigma}$ itself, since the latter can declare convergence prematurely.

\begin{algorithm}[htbp]
\caption{FH-EM: maximum-likelihood estimation of $\bm{\Sigma}$}
\label{alg:fhem}
\begin{algorithmic}[1]
\Require Per-group pure-data MLEs $\tilde{\bm{\lambda}}_j$ and Fisher information matrices $\mathbf{F}_j$, $j=1,\dots,J$, obtained from Stage~1 (isotropic-prior Laplace approximation)
\Require Initial values $\mathbf{B}^{(0)}$, $\sigma^{2(0)}$; tolerance $\mathrm{tol}$; maximum iterations $T$
\Ensure $\hat{\mathbf{B}}$, $\hat\sigma^2$
\State $\bm{\Sigma}^{(0)} \gets \mathbf{B}^{(0)}(\mathbf{B}^{(0)})^{\top} + \sigma^{2(0)}\mathbf{I}_Q$
\State $l^{(0)} \gets l_{\text{FH}}(\bm{\Sigma}^{(0)})$
\For{$t = 0, 1, 2, \dots, T-1$}
    \Statex \hspace{\algorithmicindent}\textit{E-step}
    \For{$j = 1, \dots, J$}
        \State $\mathbf{V}_j \gets \left(\mathbf{F}_j + (\bm{\Sigma}^{(t)})^{-1}\right)^{-1}$
        \State $\mathbf{m}_j \gets \mathbf{V}_j \mathbf{F}_j \tilde{\bm{\lambda}}_j$
    \EndFor
    \Statex \hspace{\algorithmicindent}\textit{M-step}
    \State $\mathbf{S} \gets \dfrac{1}{J}\displaystyle\sum_{j=1}^{J}\left(\mathbf{m}_j\mathbf{m}_j^{\top} + \mathbf{V}_j\right)$
    \State $\left(\mathbf{B}^{(t+1)}, \sigma^{2(t+1)}\right) \gets \mathrm{PPCA}(\mathbf{S}, K)$ \Comment{closed form, Eqs.~\eqref{eq:tipping-ppca_b}-\eqref{eq:tipping-ppca_sigma}}
    \State $\bm{\Sigma}^{(t+1)} \gets \mathbf{B}^{(t+1)}(\mathbf{B}^{(t+1)})^{\top} + \sigma^{2(t+1)}\mathbf{I}_Q$
    \State $l^{(t+1)} \gets l_{\text{FH}}(\bm{\Sigma}^{(t+1)})$
    \If{$\left|l^{(t+1)} - l^{(t)}\right| < \mathrm{tol}$}
    \State \textbf{break}
    \EndIf
\EndFor
\State \Return $\hat{\mathbf{B}} \gets \mathbf{B}^{(t+1)}$, $\hat\sigma^2 \gets \sigma^{2(t+1)}$
\end{algorithmic}
\end{algorithm}

\subsubsection{Solutions}

We have used the term Heywood problem loosely so far. It actually refers to three distinct failure mechanisms in this thesis, not one. Table~\ref{tab:heywood-levels} separates them and lists which method addresses each one.

\begin{table}[htbp]
\centering
\begin{tabular}{p{1.6cm}p{6.3cm}p{1.6cm}p{4.3cm}}
\toprule
Level & Issue Description & Fixed by & Fixed method \\
\midrule
Issue 1 & 
The E-step prior~\eqref{eq: Overall notation structure} depends on the current $\hat{\mathbf{B}}$, causing $\hat{\bm{\lambda}}_j$ to shrink less along the directions $\hat{\mathbf{B}}$ has already identified. The M-step then enhance the self-confirming loop. &
\textsc{iso} &
Replace the E-step prior with isotropic working prior $\tau^2\mathbf{I}$. The shrinkage weight no longer depends on the direction of $\hat{\mathbf{B}}$, so the loop is broken. \\
\addlinespace
Issue 2 &
The debias estimator~\eqref{eq:debias-Sigma} subtracts one positive semi-definite matrix from another. Nothing prevents the result from leaving the positive semi-definite cone, which drives $\hat\sigma^2$ to collapse. &
\textsc{fh-em} &
Matrix $\mathbf{S}$ in M-step~\eqref{eq:fhem_mstep} is positive definite by construction, so this collapse cannot occur. \\
\addlinespace
Issue 3 &
When the true $\sigma^2$ is small or $K$ is misspecified, the likelihood itself favors a solution at or near the boundary. A property of the model and the data, not an artifact of the algorithm. &
None &
Not addressed by \textsc{iso}, debias or \textsc{fh-em}. Would require an adjusted likelihood penalty in the M-step, which is left as future work. \\
\bottomrule
\end{tabular}
\caption{Three distinct mechanisms behind the Heywood problem in this thesis, and which correction method addresses each one.}
\label{tab:heywood-levels}
\end{table}

\iffalse
\section{FH-EM Algorithm via Woodbury}

In Section~\ref{sec: EM-wb}, we introduce the EM-algorithm via Woodbury acceleration, and have proved by simulation that Woodbury formula can reduce the runtime significantly. In this section, we manage to adapt the Woodbury structure based on our three methods of optimization, and then give the estimated accuracy and runtime as the simulated results.

\paragraph*{E-step}
Both $\tilde{\bm\lambda}_j$ and $\mathbf F_j$ in Section~\ref{subsubsec:debias} are read off from the same
isotropic-prior Laplace solve that Section~\ref{para:fhem-estep} already runs every EM round with
$\bm\Sigma^{-1}$ replaced by $\tau^{-2}\mathbf{I}_Q$. The Hessian at the isotropic mode
$\hat{\bm\lambda}_{j}^{\rm iso}$ already has the diagonal-minus-low-rank form,
\begin{equation}
    \hat{\mathbf S}_{j}^{-1}=\mathbf{D}_{j}^{\rm iso}-\mathbf{P}_{j}\mathbf{M}_{j}\mathbf{P}_{j}^\top,
    \qquad
    \mathbf{D}_{j}^{\rm iso}:=\operatorname{diag}\Big(\textstyle\sum_{i\in\mathcal I_j}y_{i\cdot}\hat{\bm p}_i\Big)+\tau^{-2}\mathbf I_Q.
    \label{eq:Shat-inv-iso}
\end{equation}
Recovering $\mathbf F_j=\hat{\mathbf S}_j^{-1}-\tau^{-2}\mathbf I$ costs nothing beyond subtracting a scalar
off the diagonal,
\begin{equation}
    \mathbf F_j=\mathbf D_j^{\mathrm F}-\mathbf P_j\mathbf M_j\mathbf P_j^\top,
    \qquad
    \mathbf D_j^{\mathrm F}:=\operatorname{diag}\Big(\textstyle\sum_{i\in\mathcal I_j}M_i\hat{\bm p}_i\Big),
    \label{eq:Fj-woodbury}
\end{equation}
where $\mathbf D_j^{\mathrm F}$, $\mathbf P_j$, $\mathbf M_j$ are already materializes. Nothing here is recomputed from scratch. Substituting Equation~\eqref{eq:Fj-woodbury} and the Woodbury form of $\bm\Sigma^{-1}$ into the FH-EM posterior precision $\mathbf F_j+\bm\Sigma^{-1}$ and collecting the diagonal and rank parts gives,
\begin{equation}
    \mathbf F_j+\bm\Sigma^{-1}
    =\underbrace{\big(\mathbf D_j^{\mathrm F}+\sigma^{-2}\mathbf I_Q\big)}_{\tilde{\mathbf D}_j}
    -\underbrace{[\,\mathbf P_j\ \ \mathbf B\,]}_{\tilde{\mathbf U}_j}\,
    \underbrace{\operatorname{blkdiag}\!\big(\mathbf M_j,\ \sigma^{-4}(\mathbf I_K+\sigma^{-2}\mathbf B^\top\mathbf B)^{-1}\big)}_{\tilde{\mathbf W}_j}\,
    \tilde{\mathbf U}_j^\top.
    \label{eq:fhem-Vj-decomp}
\end{equation}
This is exactly $\mathbf D_j-\mathbf U_j\mathbf W_j\mathbf U_j^\top$ from Equation~\eqref{eq:combined},
evaluated at $\bm p_i=\hat{\bm p}_i(\hat{\bm\lambda}_j^{\rm iso})$ instead of at each group's own anisotropic-prior
mode. The FH-EM posterior precision matrix is therefore not a new object. It is the original E-step Hessian at the first-stage isotropic mode instead of at a fresh Newton solution. No new Woodbury derivation is needed, only a new point at which the existing one is evaluated.

The only genuinely new work is a single $(I_j+K)$-sized solve, since the isotropic solve only ever formed an
$I_j$-sized $\mathbf Z_j$. Applying the Woodbury identity to Equation~\eqref{eq:fhem-Vj-decomp} gives,
\begin{equation}
    \mathbf V_j=\tilde{\mathbf D}_j^{-1}+\tilde{\mathbf D}_j^{-1}\tilde{\mathbf U}_j\,
    \tilde{\mathbf Z}_j^{-1}\,\tilde{\mathbf U}_j^\top\tilde{\mathbf D}_j^{-1},
    \qquad
    \tilde{\mathbf Z}_j:=\tilde{\mathbf W}_j^{-1}-\tilde{\mathbf U}_j^\top\tilde{\mathbf D}_j^{-1}\tilde{\mathbf U}_j,
    \label{eq:fhem-Vj-woodbury}
\end{equation}
at the same $\mathcal O(Q(I_j+K)^2+(I_j+K)^3)$ cost as the base E-step. Unlike the mode-finding E-step, this is a
single closed-form solve rather than an iterated Newton search. Equation~\eqref{eq:fh-posterior} is an exact Gaussian conditional, so there is no backtracking and no convergence check.

$\mathbf m_j$ has a similar shortcut. Equation~\eqref{eq:pure-data-lambda} gives $\tilde{\bm\lambda}_j=\mathbf F_j^{-1}\hat{\mathbf S}_j^{-1}\hat{\bm\lambda}_j^{\rm iso}$. Left-multiplying by $\mathbf F_j$ cancels $\mathbf F_j^{-1}$,
\begin{equation}
    \mathbf F_j\tilde{\bm\lambda}_j=\hat{\mathbf S}_j^{-1}\hat{\bm\lambda}_j^{\rm iso}=:\mathbf b_j.
    \label{eq:fhem-bj}
\end{equation}
So $\mathbf m_j=\mathbf V_j\mathbf b_j$ never needs the pure-data MLE $\tilde{\bm\lambda}_j$ explicitly. $\mathbf b_j$ comes from applying the diagonal-minus-low-rank form~\eqref{eq:Shat-inv-iso} directly to $\hat{\bm\lambda}_j^{\rm iso}$ at $\mathcal O(QN_j)$ and no inversion. The full E-step forms $\mathbf b_j$ and factor $\tilde{\mathbf Z}_j$ in
Equation~\eqref{eq:fhem-Vj-woodbury}, then calculate $\mathbf m_j=\mathbf V_j\mathbf b_j$. We keep
$\tilde{\mathbf D}_j^{-1},\tilde{\mathbf U}_j,\tilde{\mathbf Z}_j^{-1}$ around for the M-step below.

\paragraph*{M-step}

Equations~\eqref{eq:tipping-ppca_b}--\eqref{eq:tipping-ppca_sigma} take
$\mathbf S=\frac1J\sum_j(\mathbf m_j\mathbf m_j^\top+\mathbf V_j)$ and return its leading $K$ eigen-pairs together with $\operatorname{tr}\mathbf S$. There is no way to avoid the eigen-problem itself here. What the Woodbury structure removes is the need to ever form the $Q\times Q$ matrix that gets eigen-decomposed.

Write $\mathbf M:=[\mathbf m_1,\dots,\mathbf m_J]\in\mathbb R^{Q\times J}$. Substituting
Equation~\eqref{eq:fhem-Vj-woodbury} into $\mathbf S$ splits it into three semi-positive-definite pieces,
\begin{equation}
    \mathbf S=\frac1J\mathbf M\mathbf M^\top+\bar{\mathbf D}+\bar{\mathbf C},
    \quad
    \bar{\mathbf D}:=\frac1J\sum_{j=1}^J\tilde{\mathbf D}_j^{-1},
    \quad
    \bar{\mathbf C}:=\frac1J\sum_{j=1}^J\tilde{\mathbf D}_j^{-1}\tilde{\mathbf U}_j\tilde{\mathbf Z}_j^{-1}
    \tilde{\mathbf U}_j^\top\tilde{\mathbf D}_j^{-1}.
    \label{eq:S-operator}
\end{equation}
$\bar{\mathbf D}$ is diagonal, so summing $J$ diagonal matrices costs $\mathcal O(JQ)$, which is not a matrix product. $\bar{\mathbf C}$ is a sum of $J$ rank-$(I_j+K)$ semi-positive-definite terms, building entirely from quantities the E-step already computed. None of the three pieces needs to be added into a dense $Q\times Q$ matrix. For any $\mathbf v\in\mathbb R^Q$ and Hadamard product $\odot$, we have,
\begin{equation}
    \mathbf S\mathbf v=\frac1J\mathbf M(\mathbf M^\top\mathbf v)+\bar{\mathbf D}\odot\mathbf v +\frac1J\sum_{j=1}^J\tilde{\mathbf D}_j^{-1}\tilde{\mathbf U}_j\Big(\tilde{\mathbf Z}_j^{-1}
    \big(\tilde{\mathbf U}_j^\top(\tilde{\mathbf D}_j^{-1}\odot\mathbf v)\big)\Big),
    \label{eq:Sv-matvec}
\end{equation}
This step costs $\mathcal O\big(QJ+Q\sum_j(I_j+K))\big)$, where $\mathcal O(QJ)$ cost for the $\mathbf M\mathbf M^\top$ term,
$\mathcal O(Q)$ cost for the diagonal term and $\mathcal O(Q(N_j+K))$ cost per group for the low-rank term. This matrix-vector product is exactly what a Lanczos iteration needs~\citep{Musco2015}. Running symmetric Lanczos on $\mathbf v\mapsto\mathbf S\mathbf v$ for $m$ steps, with $m$ modestly larger than $K$, gives an $m\times m$
tri-diagonal matrix whose top $K$ Ritz pairs converge to the leading eigen-pairs $(\hat{\mathbf U}_K,\hat{\bm\Lambda}_K)$
of $\mathbf S$. The convergence rate depends on the spectral gap between $\lambda_K(\mathbf S)$ and
$\lambda_{K+1}(\mathbf S)$, not on $Q$. Since $K$ is fixed and small, the total cost is
$\mathcal O\big(K\cdot (QJ+Q\sum_j(I_j+K))\big)$ against $\mathcal O(Q^3)$ for a dense eigen-decomposition of an
explicitly formed $\mathbf S$.

$\operatorname{tr}\mathbf S$ does not need the eigen-decomposition either. From Equation~\eqref{eq:fhem-Vj-woodbury} and the cyclic property of trace,
\begin{equation}
    \operatorname{tr}\mathbf S=\frac1J\sum_{j=1}^J\Big[\|\mathbf m_j\|^2+\operatorname{tr}\tilde{\mathbf D}_j^{-1}
    +\operatorname{tr}\Big(\tilde{\mathbf Z}_j^{-1}\big(\tilde{\mathbf D}_j^{-1}\tilde{\mathbf U}_j\big)^\top
    \tilde{\mathbf D}_j^{-1}\tilde{\mathbf U}_j\Big)\Big],
    \label{eq:trS-woodbury}
\end{equation}
reusing $\tilde{\mathbf D}_j^{-1}\tilde{\mathbf U}_j$ already formed for Equation~\eqref{eq:fhem-Vj-woodbury}, at
$\mathcal O((N_j+K)^2)$ per group.

With $(\hat{\mathbf U}_K,\hat{\bm\Lambda}_K)$ and $\operatorname{tr}\mathbf S$ in hand, Equations~\eqref{eq:tipping-ppca_b}-\eqref{eq:tipping-ppca_sigma} give,
\begin{equation}
    \hat\sigma^2=\frac{\operatorname{tr}\mathbf S-\operatorname{tr}\hat{\bm\Lambda}_K}{Q-K},
    \qquad
    \hat{\mathbf B}=\hat{\mathbf U}_K\big(\hat{\bm\Lambda}_K-\hat\sigma^2\mathbf I_K\big)^{1/2},
    \label{eq:fhem-wb-mstep-final}
\end{equation}
followed by the same PLT normalization as the dense algorithm. As always, $\hat\sigma^2$ is floored at the numerical lower bound whenever the smallest retained eigenvalue in $\hat{\bm\Lambda}_K$ falls below $\hat\sigma^2$.

Putting the two steps together, one full round of Woodbury-accelerated FH-EM costs,
\begin{equation}
    \mathcal O\Big(\sum_{j=1}^J\big[Q(I_j+K)^2+(I_j+K)^3\big]+K\cdot Q\big(J+\textstyle\sum_j(I_j+K)\big)\Big),
    \label{eq:fhem-wb-cost}
\end{equation}
against $\mathcal O(JQ^2+Q^3)$ for the dense computation. This is still a reduction of one full order in $Q$. The $QKJ$ term now dominates once $J$ is the large parameter.

\subsection{Algorithm}

\begin{algorithm}[htbp]
\caption{FH-EM: maximum-likelihood estimation of $\bm{\Sigma}$}
\label{alg:fhem}
\begin{algorithmic}[1]
\Require Isotropic-prior Laplace outputs $\hat{\bm\lambda}_j^{\rm iso}$, $\mathbf D_j^{\rm iso}$, $\mathbf D_j^{\mathrm F}$, $\mathbf P_j$, $\mathbf M_j$, $j=1,\dots,J$, obtained from Stage~1
\Require Initial values $\mathbf{B}^{(0)}$, $\sigma^{2(0)}$; tolerance $\mathrm{tol}$; maximum iterations $T$; Lanczos steps $m>K$
\Ensure $\hat{\mathbf{B}}$, $\hat\sigma^2$
\State $\bm{\Sigma}^{(0)} \gets \mathbf{B}^{(0)}(\mathbf{B}^{(0)})^{\top} + \sigma^{2(0)}\mathbf{I}_Q$
\State $l^{(0)} \gets l_{\text{FH}}(\bm{\Sigma}^{(0)})$
\For{$t = 0, 1, 2, \dots, T-1$}
    \Statex \hspace{\algorithmicindent}\textit{E-step}
    \For{$j = 1, \dots, J$}
        \State $\tilde{\mathbf D}_j \gets \mathbf D_j^{\mathrm F} + \big(\sigma^{2(t)}\big)^{-1}\mathbf I_Q$
        \State $\tilde{\mathbf U}_j \gets [\,\mathbf P_j\ \ \mathbf B^{(t)}\,]$
        \State $\tilde{\mathbf W}_j \gets \operatorname{blkdiag}\Big(\mathbf M_j,\ \big(\sigma^{2(t)}\big)^{-2}\big(\mathbf I_K+\big(\sigma^{2(t)}\big)^{-1}(\mathbf B^{(t)})^\top\mathbf B^{(t)}\big)^{-1}\Big)$
        \State $\mathbf b_j \gets \mathbf D_j^{\rm iso}\hat{\bm\lambda}_j^{\rm iso} - \mathbf P_j\big(\mathbf M_j(\mathbf P_j^\top\hat{\bm\lambda}_j^{\rm iso})\big)$
        \State $\tilde{\mathbf D}_j^{-1}\tilde{\mathbf U}_j \gets$ divide each column of $\tilde{\mathbf U}_j$ by $\operatorname{diag}(\tilde{\mathbf D}_j)$
        \State $\tilde{\mathbf Z}_j \gets \tilde{\mathbf W}_j^{-1} - \tilde{\mathbf U}_j^\top(\tilde{\mathbf D}_j^{-1}\tilde{\mathbf U}_j)$
        \State Factor $\tilde{\mathbf Z}_j$ (Cholesky); cache $\tilde{\mathbf Z}_j^{-1}$
        \State $\mathbf m_j \gets \tilde{\mathbf D}_j^{-1}\mathbf b_j + (\tilde{\mathbf D}_j^{-1}\tilde{\mathbf U}_j)\Big(\tilde{\mathbf Z}_j^{-1}\big((\tilde{\mathbf D}_j^{-1}\tilde{\mathbf U}_j)^\top\mathbf b_j\big)\Big)$
    \EndFor
    \Statex \hspace{\algorithmicindent}\textit{M-step}
    \State $\mathbf M \gets [\mathbf m_1,\dots,\mathbf m_J]$
    \State $\bar{\mathbf D} \gets \dfrac{1}{J}\displaystyle\sum_{j=1}^{J} \tilde{\mathbf D}_j^{-1}$
    \State operator $\mathbf S: \mathbf v \mapsto \dfrac{1}{J}\mathbf M(\mathbf M^\top\mathbf v) + \bar{\mathbf D}\odot\mathbf v + \dfrac{1}{J}\displaystyle\sum_{j=1}^{J}(\tilde{\mathbf D}_j^{-1}\tilde{\mathbf U}_j)\Big(\tilde{\mathbf Z}_j^{-1}\big((\tilde{\mathbf D}_j^{-1}\tilde{\mathbf U}_j)^\top\mathbf v\big)\Big)$
    \State run symmetric Lanczos on $\mathbf v\mapsto\mathbf S\mathbf v$ for $m$ steps to obtain top-$K$ Ritz pairs $(\hat{\mathbf U}_K,\hat{\bm\Lambda}_K)$
    \State $\operatorname{tr}\mathbf S \gets \dfrac{1}{J}\displaystyle\sum_{j=1}^{J}\Big[\|\mathbf m_j\|^2 + \operatorname{tr}\tilde{\mathbf D}_j^{-1} + \operatorname{tr}\Big(\tilde{\mathbf Z}_j^{-1}(\tilde{\mathbf D}_j^{-1}\tilde{\mathbf U}_j)^\top(\tilde{\mathbf D}_j^{-1}\tilde{\mathbf U}_j)\Big)\Big]$
    \State $\sigma^{2(t+1)} \gets \max\Big(\dfrac{\operatorname{tr}\mathbf S - \operatorname{tr}\hat{\bm\Lambda}_K}{Q-K},\ \sigma^2_{\min}\Big)$
    \State $\mathbf{B}^{(t+1)} \gets \hat{\mathbf U}_K\big(\hat{\bm\Lambda}_K - \sigma^{2(t+1)}\mathbf I_K\big)^{1/2}$
    \State apply PLT normalization to $\mathbf{B}^{(t+1)}$
    \State $\bm{\Sigma}^{(t+1)} \gets \mathbf{B}^{(t+1)}(\mathbf{B}^{(t+1)})^{\top} + \sigma^{2(t+1)}\mathbf{I}_Q$
    \State $l^{(t+1)} \gets l_{\text{FH}}(\bm{\Sigma}^{(t+1)})$
    \If{$\left|l^{(t+1)} - l^{(t)}\right| < \mathrm{tol}$}
    \State \textbf{break}
    \EndIf
\EndFor
\State \Return $\hat{\mathbf{B}} \gets \mathbf{B}^{(t+1)}$, $\hat\sigma^2 \gets \sigma^{2(t+1)}$
\end{algorithmic}
\end{algorithm}

\section{Consistency \& Asymptotics}
\label{sec:asymptotical-properties}
